\section{Dependencies, procedures and reproduction}

The constructions used are developed within this article:
the dodecaphasic register in Section~\ref{sec:k-registro}, positional
closure in Section~\ref{sec:alpha}, action in
Section~\ref{hmtII:sec:action}, and the angular pair in
Section~\ref{hmtII:sec:ii-angulos}. The constitutive response in
Section~\ref{iii:sec:constitutiva} composes these outputs while
preserving the sheet and orientation labels.

\subsection{Order of reproduction}

Reproduction distinguishes the coordinate producer from the subsequent
evaluators. The former enumerates admissible states and applies the rules
of APP, TRIT and TPK. Its regional outputs are legitimate inputs to the
subsequent operators of the same development. The vacuum evaluation
program is not presented as a substitute for that producer.

The package reading order is: kernel definitions; regional enumeration
and refinement; signed register and transformation of $K$;
positional closure of $\alpha$; action section and return; angular pair;
contractive channels; constitutive operator; cyclic moments and the
Barbero functional; dimensional sections and subsequent comparison.

Each stage preserves its domain. A regional-digit file carries the
producer's fingerprint, not a metrological attribution. The twelve-component
tuple of $K$ has its source in the register chapter. The
action coefficients, $90/120$ incidence and $12:-1$
multiplicities are read in their proofs preceding evaluation. The
local program reproduces those compositions at the declared precision.

\subsection{Three checks with different functions}

Documentary preservation checks file identity and continuity
of the copies. Mathematical proof establishes identities and limits
in their domains. Numerical execution checks a particular evaluation,
its stability and the residuals of the identities. None of these three
checks replaces the others.

The accompanying package includes a source manifest, the LaTeX texts,
the regional data used and the verification programs. The proofs
of inversion, monotonicity, moments, transport and covariance are given
in the body of the article. Documentary sources are cited for
provenance; readers can check the arguments developed here
without using a target numerical result as a generator.

The positional closure of $\alpha$ is proved in
Proposition~\ref{prop:alpha-normalizacion}. Theorem
\ref{thm:alpha-dos-publicaciones-completas} identifies its prefixes with
those of the correlated analytic representation of the same state,
through the explicit recurrence and compatible cylinders.
Constitutive inversion in
Section~\ref{iii:sec:constitutiva} uses the specified internal
normalization: changes of units are undone before that inverse is
applied. These conditions delimit the results used by the subsequent
compositions.
